\scanpage{8}
% archive foliation: 7
2. \unsure{Structures|Structure} de carré combinatoire.

Le carré peut se définir, soit (dans le contexte des polygones généraux)
par un système
\[
(1)\qquad
(\underset{\text{sommets}}{S},\ \underset{\text{arêtes}}{A},\
\underset{\text{incidences}}{R})
\]
satisfaisant les axiomes connus, soit (en utilisant l'autom. ``antipodisme''
$\underline{a}$ de $Q$) en termes de
\diagram{\begin{tikzcd}
(2) & S \arrow[d,"\text{degré 2}"] & \\
& \begin{array}{c}D\\(\simeq S/\underline{a})\end{array}
& \begin{array}{c}\text{ens des diagonales}\\\text{(de card. 2)}\end{array}
\end{tikzcd}}
comme une famille de deux ensembles (les diagonales) de cardinal 2, soit dualement
\deleted{par} \unsure{en} termes de
\margin{Ou encore comme un ens. $S$ de cardinal 4, avec une involution donnée
par \deleted{\unsure{$a$}} $x\in V^*$, i.e. une partition de type $2,2$.}
\diagram{\begin{tikzcd}
(3) & A \arrow[d,"\text{degré 2}"] & \\
& \begin{array}{c}C\\(\simeq A/\underline{a})\end{array}
& \text{ens des codiagonales}
\end{tikzcd}}
comme une \deleted{\unsure{deux}} famille de deux ens.
(les codiagonales \inserted{ou paires de côtés opposés}) de cardinal 2.
Chacune des descriptions (1) (2) (3) est valable pour décrire le carré ---
la première est \deleted{symétrique}\inserted{autoduale}, les
\deleted{deuxième}\inserted{\unsure{deux autres}} ne le sont pas, l'une privilégiant
les points de \unsure{vue} ``sommets'', l'autre les points de \unsure{vue} ``arêtes'', comme ens de base.
Car dans l'un et l'autre cas, c'est un ens. à quatre éléments avec involution
ss points
